\section{DA LEGITIMIDADE PARA REPRESENTAÇÃO}

\begin{enumerate}[resume]
\item Em caráter inicial, observa-se que o município de \nomeMunicipio{} firmou o Contrato nº \numeroContrato{} com a empresa GRID – SOLUÇÕES EM RECUPERAÇÃO DE ATIVOS, inscrita no CNPJ nº 50.610.209/0001-45. A publicidade desse instrumento ocorreu no \publicacao{} em \dataPublicacao{}. O referido contrato tem como finalidade a proteção dos interesses da Contratante, com foco em questões de natureza administrativa perante a distribuidora de energia elétrica e os órgãos reguladores competentes. As especificidades dessa atuação encontram-se detalhadas na \tituloClausula{} (Anexo II), conforme segue abaixo:
\Citacao{\clausula{}}
\item Em face do que foi apresentado, sublinha-se a validade da empresa GRID – SOLUÇÕES EM RECUPERAÇÃO DE ATIVOS, legalmente representada por sua sócia-administradora, para agir em nome da Cliente. A companhia terá a incumbência de protocolar requerimentos e defender os interesses da Cliente frente à Agência Nacional de Energia Elétrica (ANEEL) e à concessionária de energia elétrica. 
\item Desta forma, para fazer face ao que dispõe o Parecer n. 00361/2023/PFANEEL/PGF/AGU, itens 136 e 139, seguem anexos os seguintes documentos:  
\item Instrumento de Procuração que confere poderes à Sra. Ana Livia F. Dias para representar o município contratante perante a Distribuidora de Energia Elétrica e a ANEEL, junto com os demais documentos referentes à procuração: kit prefeito, contrato social da empresa contratada e documento de identificação da representante legal da empresa (Anexo I). 
\item Contrato de Prestação de Serviços (Contrato de Mandato) celebrado entre o município e a GRID – SOLUÇÕES EM RECUPERAÇÃO DE ATIVOS, junto com os demais documentos referentes ao contrato: comprovante de publicação do contrato, portaria de nomeação do signatório do contrato (Anexo II).  
\item Todos os documentos apresentados atestam a validade e a vigência da contratação da empresa peticionante.










\item A capacidade postulatória e os poderes de representação do subscritor estão devidamente comprovados pela procuração anexada (Anexo I), outorgada pelo Prefeito Municipal. Tal instrumento é plenamente suficiente para a atuação administrativa perante a distribuidora e a ANEEL, em conformidade com o art. 9º da REN nº 1.000/2021 e demais legislações aplicáveis.
\item Qualquer exigência documental que exceda a procuração, como a apresentação do contrato de honorários, não possui amparo legal e infringe as prerrogativas da advocacia, estabelecidas pela Lei nº 8.906/1994.
\item Ainda assim, com o intuito de evitar dilações indevidas e demonstrar total transparência, anexa-se voluntariamente cópia do contrato Nº \numeroContrato{} de prestação de serviços celebrado entre a municipalidade e a sociedade de advogados em \dataPublicacao{}.
\item O referido contrato, em sua \textbf{\tituloClausula{}}, confere poderes expressos e específicos para que este representante atue em nome do Município em todas as instâncias administrativas necessárias, notadamente perante a concessionária e a ANEEL, para a recuperação de valores e defesa de seus interesses.
\Citacao{\clausula{}}
\item A cláusula transcrita confere expressamente poderes de representação ao representante do escritório contratado pelo Município, nos moldes do Contrato de Mandato já juntado aos autos, o que torna indiscutível sua legitimidade para atuar em nome do ente municipal perante esta distribuidora.
\item Adicionalmente, cumpre informar que o contrato em questão foi devidamente publicado no(a) \textbf{\publicacao{}}, em obediência à legislação vigente, ato que lhe confere plena eficácia jurídica e oponibilidade contra terceiros.
\item Dessa forma, o requisito formal de publicação do referido instrumento contratual foi integralmente cumprido. 
\end{enumerate}

